\chapter{Lipoprotein(a)}

\section*{What Is This Test?}
Lipoprotein(a) --- Lp(a) --- is a low-density lipoprotein (LDL) variant in which apolipoprotein B-100 is covalently bound to an apolipoprotein(a) moiety. Lp(a) is a genetically determined cardiovascular risk marker and prothrombotic lipoprotein. Elevated Lp(a) promotes atherosclerosis (apo-B retained in vessel wall) and thrombosis (apo(a) homology to plasminogen interferes with fibrinolysis). Levels are largely set at birth, are 90\% genetically determined, and do not respond meaningfully to lifestyle interventions. Hence Lp(a) identifies patients who may require aggressive LDL lowering (statins, PCSK9 inhibitors) despite ``normal'' lipid panels.

\section*{Alternative Names}
\begin{itemize}
  \item Lipoprotein(a)
  \item Lp(a)
  \item LpA
  \item Apolipoprotein(a)
\end{itemize}
\section*{Principle}
Quantified in serum or plasma by immunoturbidimetric or immunonephelometric assays using monoclonal or polyclonal antibodies against apolipoprotein(a). Results vary between assays because of the size polymorphism of apo(a) isoforms. A standardized assay (in nmol/L) reduces inter-laboratory discrepancies; recent IFCC standardization favors nmol/L reporting. Per specimen, fasting is not strictly required, but a fasting sample collected for standard lipids is usually used. Newer PCSK9 inhibitors and apheresis reduce levels; niacin modestly lowers but not clinically meaningful.

\section*{History}
Lp(a) was discovered by Kare Berg in 1963. Though evidence of cardiovascular risk derived from numerous epidemiologic studies in the late 20th century, clinical uptake was slow. In 2018 the European Atherosclerosis Society and American Heart Association/National Lipid Association Scientific Statements recommended Lp(a) testing once in every adult's lifetime \cite{kronenberg2022lipoprotein,wilson2019use}. Phase 3 trials of Lp(a)-lowering antisense oligonucleotides (pelacarsen, olpasiran) are ongoing \cite{viney2016antisense}.

\section*{Why This Test Is Ordered}
\begin{itemize}
  \item Once-in-a-lifetime screening of all adults (EAS, ACC/AHA)
  \item Risk stratification in patients with premature atherosclerotic cardiovascular disease (ASCVD) or strong family history
  \item Patients with recurrent cardiovascular events despite optimal lipid management
  \item Family history of elevated Lp(a) (first-degree relatives) or high-risk patient groups
  \item Children of affected parents, ideally before puberty
  \item Risk assessment in familial hypercholesterolemia
\end{itemize}

\section*{Primary vs Secondary Test}
Secondary, considered additional cardiovascular risk assessment rather than universal screening, although guidelines increasingly supportive of universal adult screening.

\section*{How This Test Is Performed}
A venous blood sample (serum or EDTA plasma) drawn. The test can be added to a standard fasting lipid panel. Because Lp(a) is essentially constant over adult life, only one measurement is needed in most people.

\section*{What Preparation Is Needed}
Fasting not strictly required; consistent with standard lipid panel timing.

\section*{Contraindications and Precautions}
\begin{itemize}
  \item Assays vary in calibration; results may differ across laboratories. Same patient should be measured with the same assay if monitoring therapy. Apo(a) size heterogeneity can bias results: small isoforms under-read, large isoforms over-read depending on assay.
\end{itemize}
\section*{Risks}
None.

\section*{What Happens After the Test}
Elevated Lp(a) is a lifelong risk factor: implement aggressive LDL-C lowering (statin, ezetimibe, PCSK9 inhibitor), consider aspirin, identify and treat other risk factors (BP, diabetes, smoking). Lipoprotein apheresis reserved for very high Lp(a) with progressive ASCVD despite maximal therapy.

\section*{Understanding Results}

\subsection*{Below Normal}
Low Lp(a) is reassuring; standard risk factor-based care continues.

\subsection*{Above Normal}
\begin{itemize}
  \item Borderline moderate risk: 30--50 mg/dL (75--125 nmol/L)
  \item High cardiovascular risk: 50--100 mg/dL (125--250 nmol/L): 2-fold increased risk
  \item Very high risk: $>$ 100 mg/dL ($>$ 250 nmol/L): up to 4-fold risk of stroke
  \item Familial: Very high in children predicts early coronary disease in adulthood
  \item \textbf{Action when elevated:} Aggressive LDL-C goal $<$ 55 mg/dL in secondary prevention; higher intensity statin + ezetimibe; PCSK9 inhibitor if needed
\end{itemize}

\section*{Normal Results}
\begin{itemize}
  \item \textbf{Optimal:} $<$ 30 mg/dL ($<$ 75 nmol/L)
  \item \textbf{Elevated (EAS, ACC/AHA):} $>$ 50 mg/dL ($>$ 125 nmol/L) confers 2--3-fold cardiovascular risk \cite{mach2020esc}
  \item \textbf{Higher risk in Caucasians}: $>$ 50 mg/dL (median population approximately 10--15 mg/dL)
  \item \textbf{Higher risk groups:} African Americans (median $\sim$ 30 mg/dL); 90\% genetically determined, stable over lifetime
\end{itemize}

\section*{Follow-up Tests}
\begin{itemize}
  \item \textbf{Standard lipid panel:} Total cholesterol, LDL-C, HDL-C, triglycerides, non-HDL-C
  \item \textbf{Apolipoprotein B:} Marker of atherogenic particle number; useful in insulin resistance, metabolic syndrome
  \item \textbf{Coronary artery calcium score (CACS):} Risk refinement in asymptomatic patients
  \item \textbf{Stepped familial ASCVD testing:} Cascade family testing for siblings and children
  \item \textbf{Lp(a)-lowering therapy options (in trials)}: Pelacarsen, olpasiran (antisense/oligonucleotide)
  \item \textbf{Lipoprotein apheresis:} Refractory patients with progressive ASCVD despite maximal therapy
\end{itemize}

\section*{References}
\bibliographystyle{unsrtnat}
\bibliography{references}

\clearpage
\section*{Regulatory Status and Authorization Date}
This chapter describes a test category, not a specific commercial product. FDA, CDSCO, EU IVDR, and MHRA approval or registration dates therefore cannot be assigned without the assay manufacturer, product name, instrument, and intended use. For laboratory-developed tests, an individual agency approval date may not exist. See the Preface for the interpretation of regulatory status.

\clearpage
